% questions.tex —— 2015 年·福建卷·文科 · 题目内容片段（无答案）
% 由 tools/generate_projects.py 依来源总账自 2015/2015_数学真题_福建文_试题.md 生成，勿手改（改源 MD 后重跑）


\begin{gktip}
固定来源：普通高考/2015/2015福建文.pdf，SHA-256 见来源校对记录。原 PDF 为 2 页 A3 横向，印刷页码 945–946，含题面但不含答案页。题面已完成独立逐题核对；独立推导答案与解析另存于对应文件，非官方答案。来源发布者/出版版次仍未核实；图形和统计表只登记在图表清单，未用扫描图片替代矢量源。
\end{gktip}

\gkpart{一、选择题}

\begin{question}
\gkti{1}

若 \((1+i)+(2-3i)=a+bi\)（\(a,b\in\mathbb R\)，\(i\) 是虚数单位），则 \(a,b\) 的值分别等于（　）

(A) \(3,-2\)　(B) \(3,2\)　(C) \(3,-3\)　(D) \(-1,4\)

\end{question}
\begin{question}
\gkti{2}

若集合 \(M=\{x\mid -2\leqslant x<2\}\)，\(N=\{0,1,2\}\)，则 \(M\cap N\) 等于（　）

(A) \(\{0\}\)　(B) \(\{1\}\)　(C) \(\{0,1,2\}\)　(D) \(\{0,1\}\)

\end{question}
\begin{question}
\gkti{3}

下列函数为奇函数的是（　）

(A) \(y=\sqrt{x}\)　(B) \(y=e^x\)　(C) \(y=\cos x\)　(D) \(y=e^x-e^{-x}\)

\Needspace{0.4\textheight}
\end{question}
\begin{question}
\gkti{4}

阅读如图所示的程序框图，运行相应的程序，若输入 \(x\) 的值为 \(1\)，则输出 \(y\) 的值为（　）

(A) \(2\)　(B) \(7\)　(C) \(8\)　(D) \(128\)
\nopagebreak[4]
\begin{center}% Source: fixed 2015 Fujian arts paper, printed page 945.
% Vector redraw; the diagram records the printed branch labels and formulas.
\begin{tikzpicture}[x=11mm,y=11mm,>=Stealth,
  every node/.style={font=\fontsize{10.5pt}{12.6pt}\selectfont,inner sep=2pt},
  every path/.style={line width=.45pt}]
  \node[draw,rounded corners] (start) at (0,0) {开始};
  \node[draw,trapezium,trapezium left angle=70,trapezium right angle=110]
    (input) at (0,-1.1) {输入 \(x\)};
  \node[draw,diamond,aspect=2.1,minimum width=25mm,minimum height=10mm]
    (test) at (0,-2.5) {\(x\geqslant2?\)};
  \node[draw,rectangle,minimum width=19mm,minimum height=8mm]
    (yes) at (-1.45,-4.0) {\(y=2^x\)};
  \node[draw,rectangle,minimum width=19mm,minimum height=8mm]
    (no) at (1.45,-4.0) {\(y=9-x\)};
  \node[draw,trapezium,trapezium left angle=70,trapezium right angle=110]
    (output) at (0,-5.25) {输出 \(y\)};
  \node[draw,rounded corners] (finish) at (0,-6.35) {结束};

  \draw[->] (start) -- (input);
  \draw[->] (input) -- (test);
  \draw[->] (test.south) -- node[left,inner sep=1pt] {是} (yes.north);
  \draw[->] (test.east) -- ++(8mm,0) |- (no.north);
  \node at (1.12,-2.62) {否};
  \coordinate (merge) at (0,-4.8);
  \draw (yes.south) -- (yes.south |- merge) -- (merge);
  \draw[->] (no.south) -- (no.south |- merge) -- (merge);
  \draw[->] (merge) -- (output.north);
  \draw[->] (output) -- (finish);
\end{tikzpicture}
\end{center}
\nopagebreak[4]
\end{question}
\begin{question}
\gkti{5}

若直线 \(\dfrac{x}{a}+\dfrac{y}{b}=1\)（\(a>0,b>0\)）过点 \((1,1)\)，则 \(a+b\) 的最小值等于（　）

(A) \(2\)　(B) \(3\)　(C) \(4\)　(D) \(5\)

\end{question}
\begin{question}
\gkti{6}

若 \(\sin\alpha=-\dfrac{5}{13}\)，且 \(\alpha\) 为第四象限角，则 \(\tan\alpha\) 的值等于（　）

(A) \(\dfrac{12}{5}\)　(B) \(-\dfrac{12}{5}\)　(C) \(\dfrac{5}{12}\)　(D) \(-\dfrac{5}{12}\)

\end{question}
\begin{question}
\gkti{7}

设 \(\boldsymbol a=(1,2)\)，\(\boldsymbol b=(1,1)\)，\(\boldsymbol c=\boldsymbol a+k\boldsymbol b\)。若 \(\boldsymbol b\perp\boldsymbol c\)，则实数 \(k\) 的值等于（　）

(A) \(-\dfrac32\)　(B) \(-\dfrac53\)　(C) \(\dfrac53\)　(D) \(\dfrac32\)

\Needspace{0.4\textheight}
\end{question}
\begin{question}
\gkti{8}

\Needspace{5\baselineskip}
如图，矩形 \(ABCD\) 中，点 \(A\) 在 \(x\) 轴上，点 \(B\) 的坐标为 \((1,0)\)，且点 \(C\) 与点 \(D\) 在函数
\nopagebreak[4]
\[
f(x)=\begin{cases}
x+1,&x\geqslant0,\\
-\dfrac12x+1,&x<0
\end{cases}
\]
的图象上。若在矩形 \(ABCD\) 内随机取一点，则此点取自阴影部分的概率等于（　）

(A) \(\dfrac16\)　(B) \(\dfrac14\)　(C) \(\dfrac38\)　(D) \(\dfrac12\)
\nopagebreak[4]
\begin{center}% Source: fixed 2015 Fujian arts paper, printed page 945.
% Vector redraw from the exact function and rectangle conditions; not a raster crop.
\begin{tikzpicture}[x=16mm,y=16mm,>=Stealth,
  every node/.style={font=\fontsize{10.5pt}{12.6pt}\selectfont,inner sep=1pt},
  every path/.style={line width=.45pt}]
  % From B=(1,0), f(1)=2 and the stated rectangle, the corners are
  % A=(-2,0), B=(1,0), C=(1,2), D=(-2,2).
  \coordinate (A) at (-2,0);
  \coordinate (B) at (1,0);
  \coordinate (C) at (1,2);
  \coordinate (D) at (-2,2);
  \coordinate (V) at (0,1);
  \coordinate (O) at (0,0);
  \fill[gray!28] (D) -- (C) -- (V) -- cycle;
  \draw (A) rectangle (C);
  \draw[thick] (D) -- (V) -- (C);
  \draw[->] (-2.25,0) -- (1.45,0) node[right] {\(x\)};
  \draw[->] (0,-.25) -- (0,2.45) node[above] {\(y\)};
  \node[below left] at (A) {\(A\)};
  \node[below right] at (B) {\(B\)};
  \node[above right] at (C) {\(C\)};
  \node[above left] at (D) {\(D\)};
  \node[below left] at (O) {\(O\)};
\end{tikzpicture}
\end{center}
\nopagebreak[4]
\Needspace{0.4\textheight}
\end{question}
\begin{question}
\gkti{9}

某几何体的三视图如图所示，则该几何体的表面积等于（　）

(A) \(8+2\sqrt2\)　(B) \(11+2\sqrt2\)　(C) \(14+2\sqrt2\)　(D) \(15\)
\nopagebreak[4]
\begin{center}% Source: fixed 2015 Fujian arts paper, printed page 945.
% Orthographic views and dimensions follow the source; proportions are schematic.
\begin{tikzpicture}[x=13mm,y=13mm,>=Stealth,
  every node/.style={font=\fontsize{10.5pt}{12.6pt}\selectfont,inner sep=1pt},
  every path/.style={line width=.45pt}]
  % Front view: 2 by 2, split into two unit-width panels.
  \draw (0,0) rectangle (2,2);
  \draw (1,0) -- (1,2);
  \draw[<->] (-.22,0) -- (-.22,2) node[midway,left] {\(2\)};
  \draw[<->] (0,-.22) -- (1,-.22) node[midway,below] {\(1\)};
  \draw[<->] (1,-.22) -- (2,-.22) node[midway,below] {\(1\)};
  \node[below=7mm] at (1,0) {正视图};

  % Side view: width 1 and height 2.
  \begin{scope}[xshift=4.0cm]
    \draw (0,0) rectangle (1,2);
    \draw[<->] (0,-.22) -- (1,-.22) node[midway,below] {\(1\)};
    \node[below=7mm] at (.5,0) {侧视图};
  \end{scope}

  % Top view: trapezoid. Only the source topology is encoded; no extra dimensions.
  \begin{scope}[xshift=1.1cm,yshift=-3.6cm]
    \draw (0,0) -- (1,0) -- (2,1) -- (0,1) -- cycle;
    \node[below=7mm] at (1,0) {俯视图};
  \end{scope}
\end{tikzpicture}
\end{center}
\nopagebreak[4]
\end{question}
\begin{question}
\gkti{10}

\Needspace{5\baselineskip}
变量 \(x,y\) 满足约束条件
\nopagebreak[4]
\[
\begin{cases}
x+y\geqslant0,\\
x-2y+2\geqslant0,\\
mx-y\leqslant0,
\end{cases}
\]
若 \(z=2x-y\) 的最大值为 \(2\)，则实数 \(m\) 等于（　）

(A) \(-2\)　(B) \(-1\)　(C) \(1\)　(D) \(2\)

\end{question}
\begin{question}
\gkti{11}

已知椭圆 \(E:\dfrac{x^2}{a^2}+\dfrac{y^2}{b^2}=1\)（\(a>b>0\)）的右焦点为 \(F\)，短轴的一个端点为 \(M\)，直线 \(l:3x-4y=0\) 交椭圆 \(E\) 于 \(A,B\) 两点。若 \(|AF|+|BF|=4\)，点 \(M\) 到直线 \(l\) 的距离不小于 \(\dfrac45\)，则椭圆 \(E\) 的离心率的取值范围是（　）

(A) \((0,\dfrac{\sqrt3}{2}]\)　(B) \((0,\dfrac34]\)　(C) \([\dfrac{\sqrt3}{2},1)\)　(D) \([\dfrac34,1)\)

\end{question}
\begin{question}
\gkti{12}

“对任意 \(x\in(0,\dfrac\pi2)\)，\(k\sin x\cos x<x\)”是“\(k<1\)”的（　）

(A) 充分而不必要条件　(B) 必要而不充分条件

(C) 充分必要条件　(D) 既不充分也不必要条件

\end{question}
\gkpart{二、填空题}

\begin{question}
\gkti{13}

某校高一年级有 900 名学生，其中女生 400 名，按男女比例用分层抽样的方法，从该年级学生中抽取一个容量为 45 的样本，则应抽取的男生人数为\gkblank{}。

\end{question}
\begin{question}
\gkti{14}

若 \(\triangle ABC\) 中，\(AC=\sqrt3\)，\(A=45^\circ\)，\(C=75^\circ\)，则 \(BC=\)\gkblank{}。

\end{question}
\begin{question}
\gkti{15}

若函数 \(f(x)=2^{|x-a|}\)（\(a\in\mathbb R\)）满足 \(f(1+x)=f(1-x)\)，且 \(f(x)\) 在 \([m,+\infty)\) 上单调递增，则实数 \(m\) 的最小值等于\gkblank{}。

\end{question}
\begin{question}
\gkti{16}

若 \(a,b\) 是函数 \(f(x)=x^2-px+q\)（\(p>0,q>0\)）的两个不同的零点，且 \(a,b,-2\) 这三个数可适当排序后成等差数列，也可适当排序后成等比数列，则 \(p+q=\)\gkblank{}。

\end{question}
\gkpart{三、解答题}

\begin{question}
\gkti{17}

等差数列 \(\{a_n\}\) 中，\(a_2=4\)，\(a_4+a_7=15\)。

(1) 求数列 \(\{a_n\}\) 的通项公式；

(2) 设 \(b_n=2^{a_n-2}+n\)，求 \(b_1+b_2+b_3+\cdots+b_{10}\) 的值。

\end{question}
\begin{question}
\gkti{18}

全网传播的融合指数是衡量电视媒体在中国网民中影响力的综合指标。根据相关报道提供的全网传播 2015 年某全国性大型活动的“省级卫视新闻台”融合指数的数据，对名列前 20 名的“省级卫视新闻台”的融合指数进行分组统计，结果如表所示。

\begin{center}
\begin{tabular}{rcr}
\toprule
组号 & 分组 & 频数 \\
\midrule
1 & \([4,5)\) & 2 \\
2 & \([5,6)\) & 8 \\
3 & \([6,7)\) & 7 \\
4 & \([7,8]\) & 3 \\
\bottomrule
\end{tabular}
\end{center}

(1) 现从融合指数在 \([4,5)\) 和 \([7,8]\) 内的“省级卫视新闻台”中随机抽取 2 家进行调研，求至少有 1 家的融合指数在 \([7,8]\) 内的概率；

(2) 根据分组统计表求这 20 家“省级卫视新闻台”的融合指数的平均数。

\Needspace{0.4\textheight}
\end{question}
\begin{question}
\gkti{19}

已知点 \(F\) 为抛物线 \(E:y^2=2px\)（\(p>0\)）的焦点，点 \(A(2,m)\) 在抛物线 \(E\) 上，且 \(|AF|=3\)。

(1) 求抛物线 \(E\) 的方程；

(2) 已知点 \(G(-1,0)\)，延长 \(AF\) 交抛物线 \(E\) 于点 \(B\)，证明：以点 \(F\) 为圆心且与直线 \(GA\) 相切的圆，必与直线 \(GB\) 相切。
\nopagebreak[4]
\begin{center}% Source: fixed 2015 Fujian arts paper, printed page 946.
% From y^2=2px, A=(2,m), |AF|=3 one obtains p=2 and m^2=8.
% The source sketch places A on the upper branch, selecting m=2sqrt(2);
% hence F=(1,0) and the second intersection on the extended AF is B=(1/2,-sqrt(2)).
\begin{tikzpicture}[x=13mm,y=13mm,>=Stealth,
  every node/.style={font=\fontsize{10.5pt}{12.6pt}\selectfont,inner sep=1pt},
  every path/.style={line width=.45pt}]
  \draw[->] (-1.35,0) -- (3.05,0) node[right] {\(x\)};
  \draw[->] (0,-2.55) -- (0,3.25) node[above] {\(y\)};
  % y^2=4x, shown parametrically so the plotted curve matches the derived data.
  \draw[thick,domain=-2.35:3.0,samples=100,smooth,variable=\t]
    plot ({\t*\t/4},{\t});
  \coordinate (O) at (0,0);
  \coordinate (G) at (-1,0);
  \coordinate (F) at (1,0);
  \coordinate (A) at (2,{2*sqrt(2)});
  \coordinate (B) at (.5,{-sqrt(2)});
  \draw (-1.25,{-sqrt(2)/6}) -- (2.25,{13*sqrt(2)/6});
  \draw (-1.25,{sqrt(2)/6}) -- (1.1,{-7*sqrt(2)/5});
  \draw (.15,{-17*sqrt(2)/10}) -- (2.35,{27*sqrt(2)/10});
  \node[below left] at (O) {\(O\)};
  \node[below] at (G) {\(G\)};
  \node[below] at (F) {\(F\)};
  \node[above right,fill=white,inner sep=2pt] at (A) {\(A\)};
  \node[below right] at (B) {\(B\)};
\end{tikzpicture}
\end{center}
\nopagebreak[4]
\Needspace{0.4\textheight}
\end{question}
\begin{question}
\gkti{20}

如图，\(AB\) 是圆 \(O\) 的直径，点 \(C\) 是圆 \(O\) 上异于 \(A,B\) 的点，\(PO\) 垂直于圆 \(O\) 所在的平面，且 \(PO=OB=1\)。

(1) 若 \(D\) 为线段 \(AC\) 的中点，求证：\(AC\perp\) 平面 \(PDO\)；

(2) 求三棱锥 \(P-ABC\) 体积的最大值；

(3) 若 \(BC=\sqrt2\)，点 \(E\) 在线段 \(PB\) 上，求 \(CE+OE\) 的最小值。
\nopagebreak[4]
\begin{center}% Source: fixed 2015 Fujian arts paper, printed page 946.
% Vector redraw of the printed cone-and-circle perspective; proportions are schematic.
\begin{tikzpicture}[x=13mm,y=13mm,>=Stealth,
  every node/.style={font=\fontsize{10.5pt}{12.6pt}\selectfont,inner sep=1pt},
  every path/.style={line width=.45pt}]
  \coordinate (O) at (0,0);
  \coordinate (A) at (-3.1,0);
  \coordinate (B) at (3.1,0);
  \coordinate (P) at (0,3.2);
  \coordinate (C) at (.802339,-.627852);
  \coordinate (D) at (-1.148831,-.313926);
  \coordinate (E) at (1.55,1.6);

  % Projected circular base: hidden rear arc dashed, visible front arc solid.
  \draw[dashed,domain=0:180,samples=50,variable=\t]
    plot ({3.1*cos(\t)},{.65*sin(\t)});
  \draw[domain=180:360,samples=50,variable=\t]
    plot ({3.1*cos(\t)},{.65*sin(\t)});
  \draw[dashed] (A) -- (B);
  \draw[dashed] (O) -- (C);

  % Cone edges and the marked interior points/segments in the source diagram.
  \draw (P) -- (A) (P) -- (B) (P) -- (C) (P) -- (D);
  \draw (A) -- (C);
  \draw (C) -- (E);
  \draw[dashed] (D) -- (E);
  \draw[dashed] (O) -- (P);

  \node[left] at (A) {\(A\)};
  \node[right] at (B) {\(B\)};
  \node[below right] at (C) {\(C\)};
  \node[below] at (D) {\(D\)};
  \node[below right] at (O) {\(O\)};
  \node[above] at (P) {\(P\)};
  \node[above right] at (E) {\(E\)};
\end{tikzpicture}
\end{center}
\nopagebreak[4]
\end{question}
\begin{question}
\gkti{21}

已知函数 \(f(x)=10\sqrt3\sin\dfrac{x}{2}\cos\dfrac{x}{2}+10\cos^2\dfrac{x}{2}\)。

(1) 求函数 \(f(x)\) 的最小正周期；

(2) 将函数 \(f(x)\) 的图象向右平移 \(\dfrac\pi6\) 个单位长度，再向下平移 \(a\)（\(a>0\)）个单位长度后得到函数 \(g(x)\) 的图象，且函数 \(g(x)\) 的最大值为 2。

① 求函数 \(g(x)\) 的解析式；

② 证明：存在无穷多个互不相同的正整数 \(x_0\)，使得 \(g(x_0)>0\)。

\end{question}
\begin{question}
\gkti{22}

已知函数 \(f(x)=\ln x-\dfrac{(x-1)^2}{2}\)。

(1) 求函数 \(f(x)\) 的单调递增区间；

(2) 证明：当 \(x>1\) 时，\(f(x)<x-1\)；

(3) 确定实数 \(k\) 的所有可能取值，使得存在 \(x_0>1\)，当 \(x\in(1,x_0)\)，恒有 \(f(x)>k(x-1)\)。

\end{question}